\documentclass{situationist-preprint}

\title{The Nothing Is Not an Entity}
\msid{PHE-16-b}
\msfamily{PHE-16}
\msdiscipline{Phenomenology}
\mstype{Consequence}
\mskeywords{the nothing; anxiety; Heidegger; Carnap; nominalization; ontological commitment}

\begin{document}
\maketitle

\begin{abstract}
PHE-16-a argues that anxiety discloses the nothing as the withdrawal of beings as a whole, and that the disclosure reveals no further entity. We derive three consequences. Heidegger's sentences about the nothing can be paraphrased, without remainder of content, as sentences about the withdrawal, and the paraphrase quantifies over no object called \enquote{the nothing}. Carnap's analysis is therefore correct as a diagnosis of the entity reading and irrelevant to the withdrawal reading. And the nothing of anxiety is not the empty world of metaphysical cosmology: the question why there is something rather than nothing is not answered or posed by the disclosure.
\end{abstract}

\section{Paraphrase}

Theorem~3.3 of PHE-16-a states that what anxiety discloses, under the name \enquote{the nothing}, is the withdrawal of beings as a whole in the sense of Definition~2.2: the lapse, for all beings together, of the significance by which they ordinarily solicit and hold a subject, while they remain present. The following is then immediate.

\begin{proposition}\label{p:para}
Every sentence of the form \enquote{the nothing $\Phi$s}, in Heidegger's use, is equivalent to a sentence of the form \enquote{in anxiety, beings as a whole withdraw in way $\Phi'$}, in which no singular term refers to an object named \enquote{the nothing}.
\end{proposition}

\begin{proof}
By Theorem~3.3 of PHE-16-a, \enquote{the nothing}, in Heidegger's use, is a name for the withdrawal disclosed in anxiety. A sentence predicating something of it predicates that thing of the withdrawal. The withdrawal is characterized in Definition~2.2 of PHE-16-a by a condition on beings and their significance, so any predication of the withdrawal can be rewritten as a predication about beings as a whole and the lapse of their significance.
\end{proof}

Three of Heidegger's sentences illustrate the paraphrase \citep{Heidegger1929}. \enquote{The nothing is encountered in anxiety} becomes: in anxiety, beings as a whole are encountered as withdrawing. \enquote{The nothing is encountered at one with beings as a whole} becomes: the withdrawal is a way beings are present, not something present beside them. \enquote{The nothing itself nihilates} becomes, as PHE-16-a proposes: the withdrawal is not an act of the subject but something that happens to her situation as a whole. None of the paraphrases contains a quantifier ranging over an object distinct from beings.

\begin{corollary}\label{c:commit}
The disclosure of the nothing carries no ontological commitment beyond beings, the subject's projects and the significance that depends on them.
\end{corollary}

\begin{proof}
By Proposition~\ref{p:para}, the content of the disclosure can be stated using only terms for beings, for the subject, and for the significance beings have for her in virtue of her projects (Lemma~3.2 of PHE-16-a). A theory whose content can be so stated is committed to nothing further.
\end{proof}

\section{Carnap's analysis}

\citet{Carnap1932} argued that \enquote{the nothing} results from treating the quantifier phrase \enquote{nothing} as a name, so that sentences predicating something of it are either malformed or, properly regimented, say that no object has the predicate. Corollary~\ref{c:commit} shows what is right in the analysis: if \enquote{the nothing} were taken to name an object, the result would be the error Carnap describes. It also shows why the analysis does not touch PHE-16-a. The withdrawal reading does not regiment \enquote{the nothing} as the negation of an existential quantifier, and so does not yield the claim that no object has the predicate; it treats \enquote{the nothing} as an idiosyncratic name for a condition of beings, and the predications are about that condition. \citet{Friedman2000} shows that the dispute between Carnap and Heidegger turned less on grammar than on the authority of logic; Proposition~\ref{p:para} concedes to Carnap the logical point and denies that it settles the phenomenological one.

\begin{objection}
Heidegger denies that logic has authority over the question of the nothing. A paraphrase in logical terms is therefore a rejection of his position, not an interpretation of it.
\end{objection}

\begin{reply}
Heidegger's denial is aimed at a logic that would decide in advance that the nothing is merely the negation of all beings, and so an object of thought formed by negation \citep{Heidegger1929}. The paraphrase makes no such decision: it says what the disclosure presents, which is a condition of beings, not their negation. It is consistent with Heidegger's claim that the nothing is more original than the logical \enquote{not}.
\end{reply}

\section{Not the empty world}

\begin{proposition}\label{p:empty}
The nothing disclosed in anxiety is not the empty world, and its disclosure bears neither way on the question why there is something rather than nothing.
\end{proposition}

\begin{proof}
The empty world is a possibility in which no beings exist. The withdrawal requires that beings remain present (Definition~2.2 of PHE-16-a), so it is not the empty world nor a presentation of it. A disclosure whose content is a condition of existing beings presents no possibility in which none exist.
\end{proof}

MET-01 answers the cosmological question by arguing that the empty world is not a stable possibility. Proposition~\ref{p:empty} shows that this answer and Theorem~3.3 of PHE-16-a are independent: the one concerns the measure of possibilities without beings, the other a way in which actual beings are given.

\begin{thebibliography}{9}\small
\bibitem[Carnap(1932)]{Carnap1932} Carnap, R. (1932). \"Uberwindung der Metaphysik durch logische Analyse der Sprache. \emph{Erkenntnis}, 2, 219--241.
\bibitem[Friedman(2000)]{Friedman2000} Friedman, M. (2000). \emph{A Parting of the Ways: Carnap, Cassirer, and Heidegger}. Open Court.
\bibitem[Heidegger(1929)]{Heidegger1929} Heidegger, M. (1929). \emph{Was ist Metaphysik?} Friedrich Cohen.
\end{thebibliography}

\end{document}
